\chapter{2013年广东卷（理）}

\section{单选题}

\begin{problem}
设集合 \(M = \{x \mid x^2 + 2x = 0, x \in \R\}\)，\(N = \{x \mid x^2 - 2x = 0, x \in \R\}\)，则 \(M \cup N =\)（\quad）

\choices
    {\( \{0\} \)}
    {\( \{0, 2\} \)}
    {\(\{-2, 0\}\)}
    {\(\{-2, 0, 2\}\)}
\end{problem}

\begin{answer}
D.
\end{answer}

\begin{solution}
由 \(x^2+2x=x(x+2)=0\)，得 \(M=\{-2,0\}\)；由 \(x^2-2x=x(x-2)=0\)，得 \(N=\{0,2\}\). 因此 \(M\cup N=\{-2,0,2\}\)，故选 D.
\end{solution}

\begin{problem}
定义域为 \(\R\) 的四个函数 \(y = x^3\)，\(y = 2^x\)，\(y = x^2 + 1\)，\(y = 2\sin x\) 中，奇函数的个数是（\quad）

\choices
    {\(4\)}
    {\(3\)}
    {\(2\)}
    {\(1\)}
\end{problem}

\begin{answer}
C.
\end{answer}

\begin{solution}
对 \(y=x^3\)，有 \((-x)^3=-x^3\)，所以它是奇函数；\(y=2^x\) 满足 \(2^{-x}\neq 2^x\)，且不可能恒有 \(2^{-x}=-2^x\)，所以既不是奇函数也不是偶函数；\(y=x^2+1\) 满足 \((-x)^2+1=x^2+1\)，是偶函数；\(y=2\sin x\) 满足 \(2\sin(-x)=-2\sin x\)，是奇函数. 因此奇函数共有 \(2\) 个，故选 C.
\end{solution}

\begin{problem}
若复数 \(z\) 满足 \(\i z = 2 + 4\i\)，则在复平面内，\(z\) 对应的点的坐标是（\quad）

\choices
    {\((2,4)\)}
    {\((2,-4)\)}
    {\((4,-2)\)}
    {\((4,2)\)}
\end{problem}

\begin{answer}
C.
\end{answer}

\begin{solution}
由 \(\i z=2+4\i\)，两边同除以 \(\i\)，得 \(z=\dfrac{2+4\i}{\i}=(2+4\i)(-\i)=4-2\i\). 所以 \(z\) 在复平面内对应的点为 \((4,-2)\)，故选 C.
\end{solution}

\begin{problem}
已知离散型随机变量 \(X\) 的分布列为

\begin{center}
\begin{tabular}{c|ccc}
\(X\) & \(1\) & \(2\) & \(3\) \\
\hline
\(P\) & \(\frac{3}{5}\) & \(\frac{3}{10}\) & \(\frac{1}{10}\)
\end{tabular}
\end{center}

则 \(X\) 的数学期望 \(E(X) =\)（\quad）

\choices
    {\(\frac{3}{2}\)}
    {\(2\)}
    {\(\frac{5}{2}\)}
    {\(3\)}
\end{problem}

\begin{answer}
A.
\end{answer}

\begin{solution}
离散型随机变量的数学期望为各取值与其概率乘积之和，因此
\(E(X)=1\times\frac35+2\times\frac3{10}+3\times\frac1{10}=\frac35+\frac6{10}+\frac3{10}=\frac32\). 故选 A.
\end{solution}

\begin{problem}
某四棱台的三视图如图所示，则该四棱台的体积是（\quad）

\bitmapfigure[max width=0.62\linewidth]{img/2013/guangdong_science/q5_fig1.png}

\choices
    {\(4\)}
    {\(\frac{14}{3}\)}
    {\(\frac{16}{3}\)}
    {\(6\)}
\end{problem}

\begin{answer}
B.
\end{answer}

\begin{solution}
由三视图可知，该四棱台的上、下底面均为正方形，边长分别为 \(1\) 和 \(2\)，高为 \(2\). 所以两底面积分别为 \(1\) 和 \(4\)，四棱台的体积为
\(V=\frac{h}{3}\left(S_1+S_2+\sqrt{S_1S_2}\right)=\frac23\left(1+4+\sqrt4\right)=\frac{14}{3}\). 故选 B.
\end{solution}

\begin{problem}
设 \(m\)，\(n\) 是两条不同的直线，\(\alpha\)，\(\beta\) 是两个不同的平面，下列命题中正确的是（\quad）

\choices
    {若 \(\alpha \perp \beta, m \subset \alpha, n \subset \beta\)，则 \(m \perp n\)}
    {若 \(\alpha \parallel \beta\)，\(m \subset \alpha\)，\(n \subset \beta\)，则 \(m \parallel n\)}
    {若 \(m \perp n\)，\(m \subset \alpha\)，\(n \subset \beta\)，则 \(\alpha \perp \beta\)}
    {若 \(m \perp \alpha\)，\(m \parallel n\)，\(n \parallel \beta\)，则 \(\alpha \perp \beta\)}
\end{problem}

\begin{answer}
D.
\end{answer}

\begin{solution}
平面垂直时，分别位于两个平面内的任意两条直线不一定垂直，所以 A 不正确；平行平面内的直线方向可以不同，所以 B 不正确；两条相交直线分别属于两个平面，不能据此推出两个平面垂直，所以 C 不正确. 由 \(m\perp\alpha\) 且 \(m\parallel n\)，得 \(n\perp\alpha\)；又 \(n\parallel\beta\)，根据“若一平面经过另一平面的垂线，则两平面垂直”的判定定理，得 \(\alpha\perp\beta\). 因此 D 正确，故选 D.
\end{solution}

\begin{problem}
已知中心在原点的双曲线 \(C\) 的右焦点为 \(F(3,0)\)，离心率等于 \(\frac{3}{2}\)，则 \(C\) 的方程是（\quad）

\choices
    {\(\frac{x^2}{4} -\frac{y^2}{\sqrt{5}} = 1\)}
    {\(\frac{x^2}{4} -\frac{y^2}{5} = 1\)}
    {\(\frac{x^2}{2} -\frac{y^2}{5} = 1\)}
    {\(\frac{x^2}{2} -\frac{y^2}{\sqrt{5}} = 1\)}
\end{problem}

\begin{answer}
B.
\end{answer}

\begin{solution}
双曲线的焦半距为 \(c=3\)，离心率满足 \(e=\frac ca=\frac32\)，所以 \(a=2\). 又 \(c^2=a^2+b^2\)，得 \(b^2=3^2-2^2=5\). 由于焦点在 \(x\) 轴上，双曲线方程为 \(\frac{x^2}{4}-\frac{y^2}{5}=1\)，故选 B.
\end{solution}

\begin{problem}
设整数 \(n \geqslant 4\)，集合 \(X = \{1, 2, 3, \cdots, n\}\).
令集合 \(S = \{(x, y, z) \mid x, y, z \in X,\ \text{且三条件 } x < y < z, y < z < x, z < x < y \text{ 恰有一个成立}\}\)，若 \((x, y, z)\) 和 \((z, w, x)\) 都在 \(S\) 中，则下列选项正确的是（\quad）

\choices
    {\((y, z, w) \in S, (x, y, w) \notin S\)}
    {\((y, z, w) \in S\)，\((x, y, w) \in S\)}
    {\((y,z,w)\notin S\)，\((x,y,w)\in S\)}
    {\((y,z,w)\notin S\)，\((x,y,w)\notin S\)}
\end{problem}

\begin{answer}
B.
\end{answer}

\begin{solution}
由 \((x,y,z)\in S\)，三个数互不相等. 分三种循环顺序讨论. 若 \(x<y<z\)，由 \((z,w,x)\in S\) 可知 \(w<x\) 或 \(w>z\). 当 \(w<x<y<z\) 时，\(w<y<z\) 和 \(w<x<y\) 分别说明 \((y,z,w)\in S\) 和 \((x,y,w)\in S\)；当 \(x<y<z<w\) 时，\(y<z<w\) 和 \(x<y<w\) 分别说明这两个三元组也都属于 \(S\).

若 \(y<z<x\)，由 \((z,w,x)\in S\) 必有 \(z<w<x\). 于是 \(y<z<w<x\)，从而 \(y<z<w\) 以及 \(y<w<x\) 成立，故 \((y,z,w)\in S\) 且 \((x,y,w)\in S\). 若 \(z<x<y\)，同理由 \((z,w,x)\in S\) 得 \(z<w<x<y\)，此时 \(z<w<y\) 以及 \(w<x<y\) 成立，故两个三元组仍都属于 \(S\). 因此正确选项为 B.
\end{solution}

\section{填空题}

\begin{problem}
不等式 \(x^{2} + x - 2 < 0\) 的解集为\fillinblank{}.
\end{problem}

\begin{answer}
\((-2,1)\).
\end{answer}

\begin{solution}
因式分解得 \(x^2+x-2=(x+2)(x-1)\). 两根为 \(-2\) 和 \(1\)，二次项系数为正，所以乘积小于零时 \(x\) 在两根之间，即解集为 \((-2,1)\).
\end{solution}

\begin{problem}
若曲线 \(y = kx + \ln x\) 在点 \((1, k)\) 处的切线平行于 \(x\) 轴，则 \(k=\)\fillinblank{}.
\end{problem}

\begin{answer}
\(-1\).
\end{answer}

\begin{solution}
函数的导数为 \(y'=k+\frac1x\). 在 \(x=1\) 处切线斜率为 \(k+1\). 切线平行于 \(x\) 轴，故斜率为零，从而 \(k+1=0\)，得 \(k=-1\).
\end{solution}

\begin{problem}
执行如图所示的程序框图，若输入 \(n\) 的值为 \(4\)，则输出 \(s\) 的值为\fillinblank{}.

\texfigure[max width=0.34\linewidth]{img_repaint/2013/guangdong_science/q11_fig1.tex}
\end{problem}

\begin{answer}
\(7\).
\end{answer}

\begin{solution}
程序开始时 \(i=1\)，\(s=1\). 当 \(i\leqslant4\) 时，每次执行 \(s=s+(i-1)\)，然后令 \(i=i+1\). 四次循环中加入的数依次为 \(0\)，\(1\)，\(2\)，\(3\)，所以输出值为 \(s=1+0+1+2+3=7\).
\end{solution}

\begin{problem}
在等差数列 \(\{a_{n}\}\) 中，已知 \(a_{3} + a_{8} = 10\)，则 \(3a_{5} + a_{7} =\fillinblank{}\).
\end{problem}

\begin{answer}
\(20\).
\end{answer}

\begin{solution}
设等差数列首项为 \(a_1\)，公差为 \(d\)，则 \(a_n=a_1+(n-1)d\). 于是 \(a_3+a_8=2a_1+9d=10\)，而
\(3a_5+a_7=3(a_1+4d)+(a_1+6d)=4a_1+18d=2(2a_1+9d)=20\).
\end{solution}

\begin{problem}
给定区域 \(D: \begin{cases}
x + 4y \geqslant 4,\\
x + y \leqslant 4,\\
x \geqslant 0,
\end{cases}\) 令点集 \(T = \{(x_0, y_0) \in D \mid x_0, y_0 \in \Z,\ (x_0, y_0)\text{ 是 } z = x + y \text{ 在 } D \text{ 上取得最大值或最小值的点}\}\)，则 \(T\) 中的点共确定\fillinblank{} 条不同的直线.
\end{problem}

\begin{answer}
\(6\).
\end{answer}

\begin{solution}
由 \(x\geqslant0\)、\(x+4y\geqslant4\) 和 \(x+y\leqslant4\) 可知可行域的顶点为 \((0,1)\)、\((0,4)\) 和 \((4,0)\). 目标函数为 \(z=x+y\)，在边界 \(x+y=4\) 上取得最大值 \(4\)，因此最大值点中的整点为 \((0,4)\)、\((1,3)\)、\((2,2)\)、\((3,1)\)、\((4,0)\)，共五点且共线. 沿边界 \(x+4y=4\) 有 \(z=1+\frac34x\)，故最小值在 \((0,1)\) 处取得. 于是五个最大值点确定一条直线，最小值点分别与这五个点确定五条不同的直线，共确定 \(1+5=6\) 条不同的直线.
\end{solution}

\section{选考题：填空题}

\begin{problem}
已知曲线 \(C\) 的参数方程为 \(\begin{cases}
x = \sqrt{2} \cos t,\\
y = \sqrt{2} \sin t,
\end{cases}\) （\(t\) 为参数），\(C\) 在点 \((1,1)\) 处的切线为 \(l\)，以坐标原点为极点，\(x\) 轴的正半轴为极轴建立极坐标系，则 \(l\) 的极坐标方程为\fillinblank{}.
\end{problem}

\begin{answer}
\(\rho(\cos\theta+\sin\theta)=2\).
\end{answer}

\begin{solution}
由参数方程消去参数，得 \(x^2+y^2=2\)，所以 \(C\) 是圆心为原点、半径为 \(\sqrt2\) 的圆. 点 \((1,1)\) 处的半径方向为 \((1,1)\)，切线与该半径垂直，故切线方程为 \(x+y=2\). 在极坐标中 \(x=\rho\cos\theta\)，\(y=\rho\sin\theta\)，代入得 \(\rho(\cos\theta+\sin\theta)=2\).
\end{solution}

\begin{problem}
如图，\(AB\) 是圆 \(O\) 的直径，点 \(C\) 在圆 \(O\) 上，延长 \(BC\) 到 \(D\) 使 \(BC = CD\)，过 \(C\) 作圆 \(O\) 的切线交 \(AD\) 于 \(E\). 若 \(AB = 6\)，\(ED = 2\)，则 \(BC =\fillinblank{}\).
\bitmapfigure[max width=0.42\linewidth]{img/2013/guangdong_science/q15_fig1.png}
\end{problem}

\begin{answer}
\(2\sqrt{3}\).
\end{answer}

\begin{solution}
设 \(BC=CD=u\)，\(AC=v\). 因为 \(AB\) 是直径，所以 \(\angle ACB=90^{\circ}\)，从而 \(AC\perp BC\)，也有 \(AC\perp CD\). 由 \(BC=CD\) 得 \(AB^2=AC^2+BC^2=AC^2+CD^2=AD^2\)，所以 \(AD=AB=6\).

建立直角坐标系，取 \(C\) 为原点，\(BC\) 所在直线为 \(x\) 轴，令 \(B=(-u,0)\)，\(D=(u,0)\)，\(A=(0,v)\). 由 \(AB=6\) 得 \(u^2+v^2=36\). 圆心为 \(O=(-\frac u2,\frac v2)\)，所以过 \(C\) 的切线方程为 \(ux-vy=0\)；直线 \(AD\) 的方程为 \(vx+uy=uv\). 设切线与 \(AD\) 的交点为 \(E\)，联立两式，得
\(E=\left(\frac{uv^2}{36},\frac{u^2v}{36}\right)\). 在从 \(A\) 到 \(D\) 的线段上，点 \(E\) 的横坐标占 \(AD\) 的比例为 \(\frac{v^2}{36}\)，因此
\(ED=AD\left(1-\frac{v^2}{36}\right)=6\cdot\frac{u^2}{36}=\frac{u^2}{6}\). 由 \(ED=2\)，得 \(u^2=12\)，故 \(BC=u=2\sqrt3\).
\end{solution}

\section{解答题}

\begin{problem}
已知函数 \(f(x) = \sqrt{2}\cos\left(x - \frac{\pi}{12}\right)\)，\(x \in \R\).

\begin{enumerate}
\item 求 \(f\left(-\frac{\pi}{6}\right)\) 的值.
\item 若 \(\cos \theta = \frac{3}{5}\)，\(\theta \in \left(\frac{3\pi}{2}, 2\pi\right)\)，求 \(f\left(2\theta + \frac{\pi}{3}\right)\).
\end{enumerate}
\end{problem}

\begin{solution}
\begin{enumerate}
\item \(f\left(-\frac{\pi}{6}\right)=\sqrt2\cos\left(-\frac\pi6-\frac\pi{12}\right)=\sqrt2\cos\left(-\frac\pi4\right)=1\).
\item 因为 \(\theta\in(\frac{3\pi}{2},2\pi)\)，所以 \(\sin\theta<0\). 由 \(\cos\theta=\frac35\)，得 \(\sin\theta=-\frac45\)，从而 \(\cos2\theta=2\cos^2\theta-1=-\frac7{25}\)，\(\sin2\theta=2\sin\theta\cos\theta=-\frac{24}{25}\). 因此
\(f\left(2\theta+\frac\pi3\right)=\sqrt2\cos\left(2\theta+\frac\pi4\right)=\sqrt2\left(\cos2\theta\cos\frac\pi4-\sin2\theta\sin\frac\pi4\right)=\frac{17}{25}\).
\end{enumerate}
\end{solution}

\begin{problem}
某车间共有 \(12\) 名工人，随机抽取 \(6\) 名，他们某日加工零件个数的茎叶图如图所示，其中茎为十位数，叶为个位数.

\begin{center}
\begin{tabular}{c|l}
1 & \(7\;9\) \\
2 & \(0\;1\;5\) \\
3 & \(0\)
\end{tabular}
\end{center}

\begin{enumerate}
\item 根据茎叶图计算样本均值.
\item 日加工零件个数大于样本均值的工人为优秀工人. 根据茎叶图推断该车间 \(12\) 名工人中有几名优秀工人？
\item 从该车间 \(12\) 名工人中，任取 \(2\) 人，求恰有 \(1\) 名优秀工人的概率.
\end{enumerate}
\end{problem}

\begin{solution}
\begin{enumerate}
\item 茎叶图中的样本数据为 \(17\)，\(19\)，\(20\)，\(21\)，\(25\)，\(30\)，所以样本均值为 \(\bar x=\frac{17+19+20+21+25+30}{6}=\frac{132}{6}=22\).
\item 样本中加工零件个数大于 \(22\) 的有 \(25\) 和 \(30\)，共 \(2\) 人，占抽取人数的 \(\frac26=\frac13\). 据此推断 \(12\) 名工人中优秀工人数为 \(12\times\frac13=4\) 名.
\item 设优秀工人有 \(4\) 名，非优秀工人有 \(8\) 名. 任取 \(2\) 人的总方法数为 \(\mathrm C_{12}^{2}\)，恰有 \(1\) 名优秀工人的方法数为 \(\mathrm C_{4}^{1}\mathrm C_{8}^{1}\)，所以所求概率为 \(\frac{\mathrm C_{4}^{1}\mathrm C_{8}^{1}}{\mathrm C_{12}^{2}}=\frac{32}{66}=\frac{16}{33}\).
\end{enumerate}
\end{solution}

\begin{problem}
如图\circled{1}，在等腰直角三角形 \(ABC\) 中，\(\angle A = 90^{\circ}\)，\(BC = 6\).
\(D\)，\(E\) 分别是 \(AC\)，\(AB\) 上的点，\(CD = BE = \sqrt{2}\)，\(O\) 为 \(BC\) 的中点.
将 \(\triangle ADE\) 沿 \(DE\) 折起，得到下图\circled{2}所示的四棱锥 \(A^{\prime} - BCDE\)，其中 \(A^{\prime}O = \sqrt{3}\).

\begin{enumerate}
\item 证明：\(A^{\prime}O \perp\) 平面 \(BCDE\).
\item 求二面角 \(A^{\prime} - CD - B\) 的平面角的余弦值.
\end{enumerate}

\bitmapfigure[max width=0.66\linewidth]{img/2013/guangdong_science/q18_fig1.png}
\end{problem}

\begin{solution}
\begin{enumerate}
\item 取平面 \(BCDE\) 为 \(xOy\) 平面，令 \(O\) 为原点，\(BC\) 在 \(x\) 轴上. 由于 \(BC=6\)，可取 \(B=(3,0,0)\)，\(C=(-3,0,0)\). 在折叠前，等腰直角三角形可取 \(A=(0,-3,0)\)，因为 \(AB=AC=3\sqrt2\). 由 \(CD=BE=\sqrt2\)，得 \(D=(-2,-1,0)\)，\(E=(2,-1,0)\)；折叠沿 \(DE\) 进行不会改变这些底面点的坐标.

设 \(A^{\prime}=(u,v,h)\). 折叠保持 \(A^{\prime}D=AD=2\sqrt2\) 和 \(A^{\prime}E=AE=2\sqrt2\)，故
\((u+2)^2+(v+1)^2+h^2=(u-2)^2+(v+1)^2+h^2\)，从而 \(u=0\). 又已知 \(A^{\prime}O=\sqrt3\)，所以 \(v^2+h^2=3\). 再由 \(A^{\prime}D^2=8\)，得 \(4+(v+1)^2+h^2=8\). 与 \(v^2+h^2=3\) 相减，得 \(2v=0\)，即 \(v=0\)，于是 \(h^2=3\). 因此 \(A^{\prime}=(0,0,\sqrt3)\)，直线 \(A^{\prime}O\) 是 \(z\) 轴，故 \(A^{\prime}O\perp\) 平面 \(BCDE\).
\item 取平面 \(BCD\) 的法向量 \(\bs{n}_2=(0,0,1)\). 在平面 \(A^{\prime}CD\) 内，取
\(\overrightarrow{CD}=(1,-1,0)\)，\(\overrightarrow{CA^{\prime}}=(3,0,\sqrt3)\)，则该平面的一个法向量为
\(\bs{n}_1=\overrightarrow{CD}\times\overrightarrow{CA^{\prime}}=(-\sqrt3,-\sqrt3,3)\). 所以二面角的平面角 \(\varphi\) 的余弦的绝对值为
\(\cos\varphi=\frac{|\bs{n}_1\cdot\bs{n}_2|}{|\bs{n}_1||\bs{n}_2|}=\frac{3}{\sqrt{3+3+9}}=\frac{3}{\sqrt{15}}=\frac{\sqrt{15}}5\).
\end{enumerate}
\end{solution}

\begin{problem}
设数列 \(\{a_{n}\}\) 的前 \(n\) 项和为 \(S_{n}\)，已知 \(a_{1} = 1\)，\(\frac{2S_{n}}{n} = a_{n+1} - \frac{1}{3} n^{2} - n - \frac{2}{3}\)，\(n \in \N^{*}\).

\begin{enumerate}
\item 求 \(a_{2}\) 的值.
\item 求数列 \(\{a_{n}\}\) 的通项公式.
\item 证明：对一切正整数 \(n\)，有 \(\frac{1}{a_1} + \frac{1}{a_2} + \cdots + \frac{1}{a_n} < \frac{7}{4}\).
\end{enumerate}
\end{problem}

\begin{solution}
\begin{enumerate}
\item 取 \(n=1\)，由 \(S_1=a_1=1\) 得 \(2=a_2-\frac13-1-\frac23=a_2-2\)，所以 \(a_2=4\).
\item 将已知关系写成 \(2S_n=na_{n+1}-\frac{n(n+1)(n+2)}3\). 当 \(n\geqslant2\) 时，与 \(n-1\) 时的关系相减，利用 \(S_n-S_{n-1}=a_n\)，得
\(2a_n=na_{n+1}-(n-1)a_n-n(n+1)\)，即 \(na_{n+1}=(n+1)a_n+n(n+1)\). 因此 \(\frac{a_{n+1}}{n+1}=\frac{a_n}{n}+1\). 又 \(\frac{a_2}{2}=2=\frac{a_1}{1}+1\)，所以令 \(b_n=\frac{a_n}{n}\)，有 \(b_1=1\)，且 \(b_{n+1}=b_n+1\)，从而 \(b_n=n\)，故 \(a_n=n^2\).
\item 由上式，原不等式等价于证明 \(\sum_{k=1}^n\frac1{k^2}<\frac74\). 当 \(n=1\)，\(2\) 时，左端分别为 \(1\) 和 \(\frac54\)，结论成立. 当 \(n\geqslant3\) 时，对 \(k\geqslant3\)，有 \(\frac1{k^2}<\frac1{k(k-1)}=\frac1{k-1}-\frac1k\)，所以
\(\sum_{k=1}^n\frac1{k^2}<1+\frac14+\sum_{k=3}^n\left(\frac1{k-1}-\frac1k\right)=\frac54+\frac12-\frac1n=\frac74-\frac1n<\frac74\). 故对一切正整数 \(n\) 均有结论成立.
\end{enumerate}
\end{solution}

\begin{problem}
已知抛物线 \(C\) 的顶点为原点，其焦点 \(F(0, c)\)（\(c > 0\)）到直线 \(l: x - y - 2 = 0\) 的距离为 \(\frac{3\sqrt{2}}{2}\). 设 \(P\) 为直线 \(l\) 上的点，过点 \(P\) 作抛物线 \(C\) 的两条切线 \(PA\)，\(PB\)，其中 \(A\)，\(B\) 为切点.

\begin{enumerate}
\item 求抛物线 \(C\) 的方程.
\item 当点 \(P(x_0, y_0)\) 为直线 \(l\) 上的定点时，求直线 \(AB\) 的方程.
\item 当点 \(P\) 在直线 \(l\) 上移动时，求 \(|AF| \cdot |BF|\) 的最小值.
\end{enumerate}
\end{problem}

\begin{solution}
\begin{enumerate}
\item 顶点在原点且焦点为 \(F(0,c)\) 的抛物线方程为 \(x^2=4cy\). 由点到直线的距离公式，得 \(\frac{c+2}{\sqrt2}=\frac{3\sqrt2}{2}\)，所以 \(c+2=3\)，即 \(c=1\). 故抛物线方程为 \(x^2=4y\).
\item 设抛物线上切点的参数分别为 \(u\)，\(v\)，则 \(A=(2u,u^2)\)，\(B=(2v,v^2)\). 抛物线在参数为 \(u\) 的点处的切线为 \(y=ux-u^2\). 因为该切线过 \(P(x_0,y_0)\)，所以 \(u\)，\(v\) 是方程 \(t^2-x_0t+y_0=0\) 的两个根，从而 \(u+v=x_0\)，\(uv=y_0\). 又因为 \(P\) 在 \(l\) 上，\(y_0=x_0-2\)，故该方程的判别式为 \(x_0^2-4y_0=(x_0-2)^2+4>0\)，确有两个不同的切点. 直线 \(AB\) 的斜率为 \(\frac{v^2-u^2}{2v-2u}=\frac{u+v}{2}=\frac{x_0}{2}\)，代入点 \(A\) 得截距 \(u^2-x_0u=-uv=-y_0\). 因此 \(AB\) 的方程为 \(y=\frac{x_0}{2}x-y_0\).
\item 由焦点 \(F(0,1)\) 及切点坐标，得 \(|AF|=\sqrt{(2u)^2+(u^2-1)^2}=u^2+1\)，同理 \(|BF|=v^2+1\). 因此
\(|AF||BF|=(u^2+1)(v^2+1)=u^2v^2+u^2+v^2+1=(uv)^2+(u+v)^2-2uv+1=x_0^2+(y_0-1)^2\). 这正是点 \(P(x_0,y_0)\) 到焦点 \(F(0,1)\) 的距离的平方. 由于 \(P\) 在直线 \(l\) 上，该距离的最小值为 \(\frac{|0-1-2|}{\sqrt{1^2+(-1)^2}}=\frac3{\sqrt2}\)，故 \(|AF||BF|\) 的最小值为 \(\left(\frac3{\sqrt2}\right)^2=\frac92\).
\end{enumerate}
\end{solution}

\begin{problem}
设函数 \(f(x) = (x - 1)\e^{x} - kx^{2} \) （\(k \in \R\)）.

\begin{enumerate}
\item 当 \(k = 1\) 时，求函数 \(f(x)\) 的单调区间.
\item 当 \(k \in \left(\frac{1}{2}, 1\right]\) 时，求函数 \(f(x)\) 在 \([0, k]\) 上的最大值 \(M\).
\end{enumerate}
\end{problem}

\begin{solution}
\begin{enumerate}
\item 当 \(k=1\) 时，\(f'(x)=x\e^x-2x=x(\e^x-2)\). 临界点为 \(x=0\) 和 \(x=\ln2\). 当 \(x<0\) 时，\(x<0\) 且 \(\e^x-2<0\)，故 \(f'(x)>0\)；当 \(0<x<\ln2\) 时，\(f'(x)<0\)；当 \(x>\ln2\) 时，\(f'(x)>0\). 所以函数的递增区间为 \((-\infty,0)\) 和 \((\ln2,+\infty)\)，递减区间为 \((0,\ln2)\).
\item 对任意给定的 \(k\in(\frac12,1]\)，有 \(f'(x)=x(\e^x-2k)\). 因为 \(2k>1\)，且由 \(\e^k\geqslant1+k\geqslant2k\)（\(k=1\) 时仍有 \(\e>2\)），可知 \(0<\ln(2k)<k\). 因此在 \([0,k]\) 上，函数先减后增，最大值只能在 \(x=0\) 或 \(x=k\) 处取得. 又 \(f(0)=-1\)，\(f(k)=(k-1)\e^k-k^3\). 为比较两端点，令 \(g(t)=1+t+t^2-\e^t\). 在 \(0\leqslant t\leqslant1\) 上，\(g'(t)=1+2t-\e^t\)，且 \(g''(t)=2-\e^t\)，所以 \(g'\) 先增后减；又 \(g'(0)=0\)，\(g'(1)=3-\e>0\)，从而 \(g'(t)\geqslant0\)，即 \(\e^t\leqslant1+t+t^2\). 于是
\(f(k)-f(0)=1-(1-k)\e^k-k^3\geqslant1-(1-k)(1+k+k^2)-k^3=0\). 所以最大值在 \(x=k\) 处取得，且 \(M=(k-1)\e^k-k^3\).
\end{enumerate}
\end{solution}
